\documentclass[11pt]{article}
\usepackage[margin=1in]{geometry}
\usepackage{amsmath,amssymb,booktabs}
\usepackage{microtype}
\usepackage{hyperref}
\usepackage{listings}
\usepackage{xcolor}

\hypersetup{colorlinks=true, linkcolor=blue, urlcolor=blue, citecolor=blue}
\definecolor{codegray}{gray}{0.95}
\lstset{language=Python,basicstyle=\small\ttfamily,backgroundcolor=\color{codegray},frame=single,breaklines=true,columns=fullflexible,showstringspaces=false}

\title{Rotary Positional Embeddings from First Principles\\
\large Giving Causal Self-Attention Information About Order}
\author{Learning Machine Learning: DeepSeek-V3 Foundations}
\date{November 2026}

\begin{document}
\maketitle

\begin{abstract}
This paper develops a small executable implementation of Rotary Positional
Embeddings (RoPE). Causal self-attention compares content vectors, but content
alone does not tell the model whether a vector occurred first or last. RoPE
introduces position by rotating pairs of query and key features by
position-dependent angles. The rotation preserves vector norms and supplies
positional information before the query-key dot product.
\end{abstract}

\tableofcontents
\newpage

\section{Introduction}
Attention can compare two hidden vectors regardless of where they occur in a
sequence. If the same vectors are moved to different positions, however, a
content-only attention mechanism has no direct representation of that change.
Autoregressive language models need order information.

Rotary Positional Embeddings address this by rotating pairs of features in the
query and key vectors. The rotation angle depends on token position and on the
feature pair's frequency. Values are not rotated because they carry the
content retrieved after query-key matching.

\section{Problem formulation}
Let
\[
X\in\mathbb{R}^{B\times S\times D}
\]
be a batch of hidden vectors. RoPE requires $D$ to be even so features can be
grouped into pairs:
\[
[x_0,x_1],[x_2,x_3],\ldots,[x_{D-2},x_{D-1}].
\]

For position $p$ and feature pair $r$, use the rotation
\[
\begin{bmatrix}
\cos\theta_{p,r} & -\sin\theta_{p,r}\\
\sin\theta_{p,r} & \cos\theta_{p,r}
\end{bmatrix}
\begin{bmatrix}x_{2r}\\x_{2r+1}\end{bmatrix}.
\]

\section{Method}
\subsection{Inverse frequencies and angles}
With base $10000$, define
\[
\omega_r=10000^{-2r/D},\qquad
\theta_{p,r}=p\omega_r.
\]
The first pair changes fastest and later pairs change more slowly.

\begin{lstlisting}
pair_indices = torch.arange(0, hidden_size, 2, device=x.device)
inverse_frequencies = base ** (-pair_indices.float() / hidden_size)
angles = position_ids[:, None].float() * inverse_frequencies[None, :]
\end{lstlisting}

\subsection{Pair rotation}
Define
\[
\operatorname{rotate\_half}(x)=[-x_1,x_0,-x_3,x_2,\ldots].
\]
Repeating the cosine and sine across each pair gives the compact equation
\[
\operatorname{RoPE}(x,p)=x\odot\cos(\theta_p)
+\operatorname{rotate\_half}(x)\odot\sin(\theta_p).
\]

\begin{lstlisting}
def rotate_half(x):
    even = x[..., ::2]
    odd = x[..., 1::2]
    return torch.stack((-odd, even), dim=-1).flatten(-2)
\end{lstlisting}

At position $p=0$, every angle is zero, so $\cos(0)=1$ and $\sin(0)=0$;
the input is unchanged.

\subsection{Norm preservation}
Each pair is multiplied by a two-dimensional orthogonal rotation matrix, so
\[
\left\|\operatorname{RoPE}(x,p)\right\|_2=\|x\|_2.
\]
RoPE changes direction while preserving vector magnitude up to floating-point
precision.

\subsection{RoPE in attention}
If $Q$ and $K$ are projected queries and keys, apply RoPE before the score
calculation:
\[
Q'=\operatorname{RoPE}(Q),\qquad K'=\operatorname{RoPE}(K),
\]
\[
L=Q'(K')^\mathsf{T},\qquad
O=\operatorname{softmax}(L+M)V.
\]
Values remain unrotated because they carry the content retrieved by attention.

\section{Reference implementation}
The implementation accepts $x$ with shape $B\times S\times D$ and shared
position IDs with length $S$.

\begin{lstlisting}
import torch


def rotate_half(x):
    even = x[..., ::2]
    odd = x[..., 1::2]
    return torch.stack((-odd, even), dim=-1).flatten(-2)


def rotary_embedding(x, position_ids=None, base=10000.0):
    batch_size, sequence_length, hidden_size = x.shape
    if hidden_size % 2 != 0:
        raise ValueError("hidden_size must be even for RoPE")

    if position_ids is None:
        position_ids = torch.arange(sequence_length, device=x.device)
    position_ids = position_ids.to(device=x.device)

    pair_indices = torch.arange(0, hidden_size, 2, device=x.device)
    inverse_frequencies = base ** (-pair_indices.float() / hidden_size)
    angles = position_ids[:, None].float() * inverse_frequencies[None, :]
    cos = angles.cos().repeat_interleave(2, dim=-1)
    sin = angles.sin().repeat_interleave(2, dim=-1)
    cos = cos.unsqueeze(0).expand(batch_size, -1, -1)
    sin = sin.unsqueeze(0).expand(batch_size, -1, -1)

    return x * cos + rotate_half(x) * sin


def main():
    torch.manual_seed(0)
    hidden_states = torch.randn(2, 4, 8)
    rotated = rotary_embedding(hidden_states)
    print("input shape:", hidden_states.shape)
    print("rotated shape:", rotated.shape)
    assert rotated.shape == hidden_states.shape

    before = torch.linalg.vector_norm(hidden_states, dim=-1)
    after = torch.linalg.vector_norm(rotated, dim=-1)
    assert torch.allclose(before, after, atol=1e-6)


if __name__ == "__main__":
    main()
\end{lstlisting}

The shape is unchanged:
\[
B\times S\times D\longrightarrow B\times S\times D.
\]

\section{A small tensor illustration}
For $x=[1,0]$, position zero gives $[1,0]$. At a position with angle
$\theta$, the output is
\[
[1,0]\begin{bmatrix}\cos\theta&-\sin\theta\\
\sin\theta&\cos\theta\end{bmatrix}
=[\cos\theta,\sin\theta].
\]
Its norm remains one because
$\sqrt{\cos^2\theta+\sin^2\theta}=1$.

\section{Executable properties}
\begin{lstlisting}
x = torch.tensor([[[1.0, 0.0, 1.0, 0.0]]])
rotated_zero = rotary_embedding(x, position_ids=torch.tensor([0]))
rotated_one = rotary_embedding(x, position_ids=torch.tensor([1]))

assert torch.allclose(rotated_zero, x)
assert not torch.allclose(rotated_zero, rotated_one)
\end{lstlisting}

The first assertion verifies the zero-angle identity. The second verifies that
position changes the representation.

\section{Discussion}
RoPE supplies order information while preserving dimensionality and norm. In
a full Transformer it is normally applied to query and key representations
inside attention. Exact tensor layouts and cache details can differ, but the
educational core remains a position-dependent rotation of feature pairs.

This paper does not reproduce DeepSeek-V3's complete attention implementation;
it isolates the positional operation needed before studying larger components.

\section{Conclusion}
Rotary Positional Embeddings turn token position into a rotation of paired
features. Applying RoPE to queries and keys before their dot product gives
causal self-attention access to both content and order, extending the Lesson 4
attention mechanism toward a more complete Transformer block.

\begin{thebibliography}{9}
\bibitem{su2021roformer}
Jianlin Su et al., ``RoFormer: Enhanced Transformer with Rotary Position
Embedding,'' 2021.\url{https://arxiv.org/abs/2104.09864}

\bibitem{deepseek}
DeepSeek-AI, ``DeepSeek-V3,'' source repository, 2024.
\url{https://github.com/deepseek-ai/DeepSeek-V3}
\end{thebibliography}

\end{document}
